\section{Discussion}

The discussion will evaluate whether temporal recovery is currently a coherent measurement family or a label applied to heterogeneous post-event outcomes. Particular attention will be given to construct drift between reaction, regulation, and recovery; whether within-person reference states are reproducible; and whether censoring and technical missingness are handled in ways that preserve interpretation.

A key boundary is that the review cannot establish clinical validity of a specific recovery endpoint merely because similar temporal constructs appear across studies. Convergence in measurement architecture can justify design choices, but validity remains population-, task-, and purpose-specific.

A focused inspection of two full-text task-interruption studies highlights the boundary between temporal trajectory, first-action latency, and error-defined restoration of sequential control \parencite{altmann2007timecourse,altmann2014momentary}. The 2007 work models a serial-position RT curve, whereas the 2014 work shows that a selective increase in sequence errors can coexist with a weak or absent RT cost. The latter's internal-consistency coefficients describe sequence-error proportions rather than reliability of time to stable recovery. Neither article supplies an individualized baseline-tolerance-stability-censoring definition of \emph{time to recovery}, and neither supports a clinical inference for children with ADHD. These primary-source observations are methodological examples; final review inclusion and frequencies await independent screening and a reproducible PRISMA flow.
